\documentclass[sigconf,natbib=false]{article}

% ── Page geometry (SIGMOD two-column, 8.5×11) ──────────────────────────────
\usepackage[letterpaper,
            top=1in, bottom=1in,
            left=0.75in, right=0.75in,
            columnsep=0.25in]{geometry}
\usepackage{multicol}
\setlength{\columnsep}{0.25in}

% ── Core packages ──────────────────────────────────────────────────────────
\usepackage{amsmath,amssymb,amsthm}
\usepackage{microtype}
\usepackage{booktabs}
\usepackage{tabularx}
\usepackage{multirow}
\usepackage{graphicx}
\usepackage{subcaption}
\usepackage{caption}
\usepackage{xcolor}
\usepackage{enumitem}
\usepackage{hyperref}
\usepackage{cleveref}
\usepackage{url}
\usepackage{parskip}
\emergencystretch=2em
\usepackage{titlesec}
\usepackage{fancyhdr}
\usepackage{abstract}
\usepackage{float}

% ── Algorithm packages ─────────────────────────────────────────────────────
\usepackage{algorithm}
\usepackage{algpseudocode}

% ── Typography ─────────────────────────────────────────────────────────────
\usepackage{times}
\usepackage{tikz}
\usetikzlibrary{shapes.geometric,arrows.meta,positioning,calc}

% ── Color palette ─────────────────────────────────────────────────────────
\definecolor{sigblue}{RGB}{26,26,46}
\definecolor{ruleblue}{RGB}{44,62,80}

% ── Section heading style ─────────────────────────────────────────────────
\titleformat{\section}[block]{\normalfont\large\bfseries\color{sigblue}}{\thesection.}{0.5em}{}[\vspace{-2pt}\rule{\columnwidth}{0.4pt}]
\titleformat{\subsection}[block]{\normalfont\normalsize\bfseries}{\thesubsection}{0.5em}{}
\titleformat{\subsubsection}[runin]{\normalfont\normalsize\bfseries}{}{0pt}{}[.\enspace]
\titlespacing*{\section}{0pt}{10pt}{4pt}
\titlespacing*{\subsection}{0pt}{8pt}{2pt}

% ── Abstract box style ─────────────────────────────────────────────────────
\renewenvironment{abstract}
  {\small\quotation\noindent\textbf{Abstract.}\enspace}
  {\endquotation}

% ── Caption style ──────────────────────────────────────────────────────────
\captionsetup{font=small,labelfont=bf,skip=4pt}
\captionsetup[table]{position=top}

% ── Hyperref ───────────────────────────────────────────────────────────────
\hypersetup{
  colorlinks=true,
  linkcolor=ruleblue,
  citecolor=ruleblue,
  urlcolor=ruleblue
}

% ── Bibliography ───────────────────────────────────────────────────────────
\usepackage[numbers,sort&compress]{natbib}
\graphicspath{{./}{figures/}}
\bibliographystyle{abbrvnat}

% ── Convenience macros ─────────────────────────────────────────────────────
\newcommand{\postgres}{PostgreSQL}
\newcommand{\arm}[1]{\texttt{#1}}
\newcommand{\feat}[1]{\textit{#1}}
\newcommand{\sys}{\textsc{RiskSteer}}
\newcommand{\eg}{\textit{e.g.},}
\newcommand{\ie}{\textit{i.e.},}
\newcommand{\etal}{\textit{et al.}}
\newcommand{\regret}{\mathcal{R}}
\newcommand{\threshold}{\tau}
\newcommand{\policy}{\pi}
\newcommand{\oracle}{\mathrm{oracle}}

% ── Theorem environments ───────────────────────────────────────────────────
\newtheorem{definition}{Definition}
\newtheorem{proposition}{Proposition}

% ── Header / footer ────────────────────────────────────────────────────────
\pagestyle{fancy}
\fancyhf{}
\renewcommand{\headrulewidth}{0pt}
% header removed
\fancyfoot[C]{\small\thepage}

% ═══════════════════════════════════════════════════════════════════════════
\begin{document}
% ═══════════════════════════════════════════════════════════════════════════

\twocolumn[{%
\begin{center}
  {\LARGE\bfseries Lightweight Risk-Controlled Planner Steering in \postgres\par}
  \vspace{10pt}
  {\large James Daramola \quad and \quad Priya Deshpande\par}
  \vspace{4pt}
  {\normalsize Department of Electrical and Computer Engineering, Marquette University\par}
  {\normalsize \texttt{james.daramola@marquette.edu}\quad
               \texttt{priya.deshpande@marquette.edu}\par}
  \vspace{12pt}
\end{center}

\begin{abstract}
Cost-based query optimizers rely on cardinality estimates that are frequently inaccurate for
multi-join workloads, causing plan selections that are orders of magnitude suboptimal in
execution time. While learned query optimization has demonstrated that machine learning can
identify and correct such failures, existing approaches depend on deep plan representations,
online exploration, or invasive changes to optimizer internals, requirements that complicate
practical deployment. We present \sys{}, a lightweight, non-intrusive planner steering
framework for \postgres{} that selects among four interpretable planner configurations using
only pre-execution \texttt{EXPLAIN} features. We introduce a \emph{risk-controlled} policy
that falls back to the default planner when the predicted benefit of steering falls below a
tunable threshold $\threshold$, directly targeting the no-regression requirement central to
production deployment.

We conduct a two-scale empirical evaluation. On the canonical 113-query Join Order Benchmark
(JOB), executed under a randomized warm-cache protocol (1{,}356 total runs), no single planner
configuration dominates: \arm{no\_hashjoin} is optimal for 34.5\% of queries,
\arm{no\_mergejoin} for 31.9\%, the default for 25.7\%, and \arm{no\_nestloop} for 8.0\%.
The default planner's tail behavior reaches a mean regret of 362.8\,ms and a p99 of
7{,}222\,ms. A Random Forest classifier trained on pre-execution features achieves 93.4\%
oracle arm prediction accuracy against a post-execution upper bound of 95.6\%, confirming
that nearly all discriminative signal is available before query execution. The dominant
predictive feature is \texttt{root\_total\_cost} (45.8\% importance), with structural plan
features accounting for the remainder.

To address the sparsity of the 113-query baseline, we scale the workload to 904 query
variants via template-based expansion, producing 3{,}616 execution samples. Across five
model families, Random Forest and Extra Trees achieve the highest accuracy at 67.96\% on the
held-out test set, substantially outperforming logistic regression (40.88\%) and confirming
that the arm-selection boundary is nonlinear. On the scaled workload, the learned policy
reduces total execution time by 37.3\%, nearly five times the 7.8\% achieved by the
best-fixed arm, and comes within 0.3\% of oracle performance. These results demonstrate
that meaningful tail latency reduction is achievable with a four-arm action space and
interpretable pre-execution features, without modifying the optimizer or relying on
heavyweight models.
\end{abstract}
\vspace{8pt}
}]

% ── CCS concepts (abbreviated) ─────────────────────────────────────────────
\noindent\textbf{CCS Concepts:} Information systems $\to$ \textit{Query optimization}; Computing methodologies $\to$ \textit{Supervised learning}.\par
\noindent\textbf{Keywords:} Query optimization, planner steering, tail latency, regret, PostgreSQL, risk-controlled policy, Join Order Benchmark
\par\vspace{6pt}
\rule{\columnwidth}{0.4pt}
\vspace{4pt}

% ═══════════════════════════════════════════════════════════════════════════
\section{Introduction}
% ═══════════════════════════════════════════════════════════════════════════

Cost-based query optimizers have been central to relational database systems for four
decades. Advances in dynamic programming plan enumeration, multi-dimensional histogram
statistics, and heuristic pruning have produced optimizers that reliably generate efficient
plans for simple queries. For complex analytical workloads, those involving many joins, correlated
predicates, and large intermediate result sets, these optimizers face a fundamental limitation:
cardinality estimation error.

Small errors in cardinality estimation compound multiplicatively through join trees. A
two-order-of-magnitude underestimate at a single join node can propagate upward, causing the
optimizer to commit to a join strategy that is catastrophically slow when executed against
actual data. The Join Order Benchmark (JOB)~\cite{Leis2015} was designed specifically to
expose this failure mode: its 113 queries over the IMDB dataset exercise correlated predicates
across 5 to 16 tables, and evaluations of \postgres{} on JOB routinely reveal that the default
planner selects plans that are 2--$100\times$ slower than the best available alternative for
a non-trivial fraction of queries.

Crucially, these failures are \emph{non-uniformly distributed} across the query set. While
the average-case behavior of \postgres{}'s planner is adequate for most queries, a small
number of queries account for the bulk of avoidable latency. This is the tail latency problem:
p95 and p99 execution times can be dominated by the optimizer's worst-case mistakes, and for
production workloads with latency service-level objectives, these outliers are
disproportionately harmful.

Learned query optimization (LQO) has emerged as the primary research response to this
failure mode. Systems such as Bao~\cite{Marcus2021Bao}, FASTgres~\cite{Habich2023},
AutoSteer~\cite{Anneser2023}, and Lero~\cite{Zhu2023Lero} demonstrate that machine learning
can identify planning failures and steer the optimizer away from them. However, these systems
make design choices that limit practical adoption. Bao requires a tree-convolutional neural
network over plan trees and online Thompson sampling exploration. Lero requires executing
multiple candidate plans per query during training. AutoSteer requires automated hint-set
discovery across multiple execution trials. The Microsoft QO-Advisor~\cite{Ding2022} targets
production deployment but operates over industrial-scale workloads with proprietary plan
representations. End-to-end systems such as Neo~\cite{Marcus2019Neo} and
Balsa~\cite{Yang2022Balsa} require reinforcement learning over full plan spaces, with
convergence times of hours to days.

\paragraph{This work.} We investigate a different point in the design space: \emph{how much
learning is necessary to reduce tail latency in \postgres{} without modifying the optimizer
or relying on heavyweight models?} We restrict the planner action space to four
configurations: the default planner and three single-operator disablement settings
(\arm{no\_hashjoin}, \arm{no\_mergejoin}, \arm{no\_nestloop}), and we train a Random Forest
classifier using only features available from \texttt{EXPLAIN (FORMAT JSON)} output before
execution. We introduce a risk-controlled steering policy that falls back to the default
planner when predicted improvement does not exceed a safety threshold $\threshold$, directly
addressing the regression avoidance requirement identified in production deployments.

Our contributions are:

\begin{itemize}[noitemsep,topsep=2pt,leftmargin=*]
  \item A \textbf{minimal action space} formalization of planner steering as a contextual
        arm selection problem with explicit regret definitions (\S\ref{sec:method}).

  \item A \textbf{risk-controlled steering policy} with a tunable safety threshold that
        controls the precision-recall tradeoff between performance gains and regression risk
        (\S\ref{sec:method:policy}).

  \item A \textbf{reproducible experimental protocol} for JOB on \postgres{}: randomized
        execution order, warm-cache regime, structured JSON plan logging, and regret-centric
        evaluation including p95/p99 tail metrics (\S\ref{sec:setup}).

  \item \textbf{Empirical characterization} of the default planner's tail behavior and the
        oracle arm distribution across 113 JOB queries (\S\ref{sec:results}), including the
        finding that no single arm dominates, which motivates learned arm selection.

  \item A \textbf{pre-execution Random Forest classifier} achieving 93.4\% oracle arm
        prediction accuracy, compared to a post-execution upper bound of 95.6\%, demonstrating
        that most steering signal is available before query execution (\S\ref{sec:ml}).
\end{itemize}

% ═══════════════════════════════════════════════════════════════════════════
\section{Background}
\label{sec:background}
% ═══════════════════════════════════════════════════════════════════════════

\subsection{PostgreSQL Query Optimization}

\postgres{} employs a cost-based optimizer that uses dynamic programming to enumerate join
orders for queries with few tables ($\leq$ \texttt{geqo\_threshold}, default 12) and a
Genetic Query Optimizer (GEQO) for larger join graphs~\cite{PGDocs}. Plan cost estimates
derive from column statistics maintained by \texttt{ANALYZE}, which samples table rows to
construct histograms and most-common-value lists. The optimizer selects among three join
strategies, namely hash join, merge join, and nested loop, based on estimated cost.

When cardinality estimates are inaccurate, the cost model may systematically prefer a join
strategy that is correct for estimated cardinalities but wrong for actual ones. This is
particularly common for join predicates involving correlated attributes, where the independent
attribute assumption underlying most histogram-based estimators breaks down.

\postgres{} exposes session-level GUC parameters that allow individual join operators to be
disabled for plan selection: \texttt{enable\_hashjoin}, \texttt{enable\_mergejoin}, and
\texttt{enable\_nestloop}. Setting any of these to \texttt{off} instructs the planner to
prefer alternative join strategies. Importantly, the \postgres{} documentation notes that
\texttt{enable\_nestloop = off} \emph{discourages} rather than fully suppresses nested loop
joins; this semantic distinction affects result interpretation and must be acknowledged in
experimental design~\cite{PGDocs}.

\subsection{The Join Order Benchmark}

The Join Order Benchmark (JOB)~\cite{Leis2015} consists of 113 queries across 33 templates
over the IMDB dataset. Queries involve between 4 and 17 table joins with correlated string
predicates on movie metadata. JOB was designed explicitly to stress join ordering and
cardinality estimation: prior work demonstrates that \postgres{} produces suboptimal plans on
a substantial fraction of JOB queries due to systematic cardinality underestimation on
correlated predicates.

FASTgres~\cite{Habich2023} observes that JOB's 113 queries are distributed across 33 structural
contexts (query templates), making per-context supervised learning challenging due to data
sparsity. This motivates template-aware evaluation splits in our experimental protocol
(\S\ref{sec:setup}). Lero~\cite{Zhu2023Lero} addresses the small dataset size by generating
1{,}000 JOB-derived queries from templates for large-scale evaluation, a direction we
identify as future work.

\subsection{Planner Configuration as Arm Selection}

We model the planner configuration selection problem as a contextual multi-armed bandit. Each
planner configuration is an \emph{arm}; the agent observes query features and selects an arm
to minimize execution time. The oracle arm is the arm that achieves minimum execution time
for a given query. The oracle arm is not known before execution but is approximable offline from experimental
traces. Regret measures the performance gap between a chosen policy and oracle arm selection.

% ═══════════════════════════════════════════════════════════════════════════
\section{Related Work}
\label{sec:related}
% ═══════════════════════════════════════════════════════════════════════════

\subsection{Steering-Based Learned Optimizers}

\textbf{Bao}~\cite{Marcus2021Bao} introduced the canonical formulation of planner steering as
a contextual bandit over hint sets. Bao encodes query plans as tree-structured inputs to a
convolutional neural network and selects per-query hint sets using Thompson sampling. Its
empirical evaluation demonstrates large improvements on JOB and other workloads, with
particular emphasis on tail latency reduction. Bao's design requires online exploration during
deployment, creating a warmup period before the model is reliable. Our approach differs in
eliminating online exploration: we train offline on collected traces and use a safety fallback
to handle low-confidence decisions without incurring regret from exploratory execution.

\textbf{FASTgres}~\cite{Habich2023} reformulates hint set selection as supervised
classification, training context-aware models (one per structural query context). It evaluates
on JOB, Stack, and TPC-H and explicitly quantifies the challenge of JOB's sparse 33-context
structure for learning. FASTgres highlights that classification accuracy is context-dependent
and that evaluation splits must account for template structure. Our experimental protocol
directly incorporates this guidance through planned template-aware cross-validation.

\textbf{AutoSteer}~\cite{Anneser2023} generalizes steering to multiple DBMSs by automatically
discovering effective hint sets through offline execution. It evaluates on JOB, Stack, TPC-DS,
and production workloads using a plan-tree convolutional model (TCNN). AutoSteer's hint
discovery process requires exhaustive cross-arm execution, which we avoid by fixing a small,
predetermined action space, trading hint-set coverage for deployment simplicity and reduced
data collection cost.

\textbf{Lero}~\cite{Zhu2023Lero} argues that predicting absolute latency is inherently
unstable and proposes learning-to-rank plans using pairwise comparisons. Lero expands JOB to
1{,}000 queries generated from templates to enable robust ML evaluation. Its pairwise ranking
framing aligns with our regret metric philosophy: both treat relative performance rather than
absolute latency as the learning signal. Our classification approach is simpler but achieves
high accuracy on the four-arm problem without requiring pairwise plan execution.

\subsection{Production-Scale Steering}

\textbf{Microsoft QO-Advisor}~\cite{Ding2022} describes the deployment of a steered optimizer
in production at Microsoft. A key contribution is its explicit treatment of regression
avoidance: the system uses a validation stage that rejects plan changes predicted to cause
regressions above an acceptable threshold. Our risk-controlled policy is directly motivated by
this framing: we operationalize a similar accept/reject mechanism via a predicted-improvement
threshold $\threshold$, making regression avoidance a first-class design constraint rather
than a post-hoc evaluation metric.

An earlier steering system~\cite{Marcus2021Steering} applies Bao-style hint selection to
industrial optimizer workloads, demonstrating the practical challenges of large search spaces
and operational constraints. This work reinforces the value of minimal action spaces: reducing
the arm count simplifies both training and deployment, at the cost of hint-set expressiveness.

\subsection{End-to-End Learned Optimizers}

\textbf{Neo}~\cite{Marcus2019Neo} was the first system to apply deep reinforcement learning
to query plan generation, but it suffers from long convergence times and instability on
unseen queries. \textbf{Balsa}~\cite{Yang2022Balsa} improves sample efficiency by introducing
simulation-to-real transfer, reducing catastrophic exploration while still requiring
substantial training infrastructure. These systems represent a different point in the
design space from steering approaches: rather than augmenting the existing optimizer, they
attempt to replace it. Deployment complexity and generalization risk make end-to-end learned
optimizers less suitable for production settings than steering methods.

\subsection{Robustness and Regression Control}

\textbf{Eraser}~\cite{Zhu2024Eraser} addresses performance regression in learned query
optimization directly, proposing mechanisms to identify and eliminate plans that cause
degradation. \textbf{Kepler}~\cite{Marcus2023Kepler} focuses on parametric query optimization
with uncertainty-aware fallbacks. Kantere and colleagues directly address regression
avoidance through a risk-aware learned cost model that formalizes robustness as a
first-class property~\cite{Kamali2024Roq}, and survey the broader landscape of
ML-based robust query optimization~\cite{Kamali2024ICDE}. Our threshold policy provides a
simpler but complementary mechanism for regression control, without requiring a probabilistic
cost model or a multi-stage validation pipeline. All of these reinforce the emerging consensus
that regression control is essential for any practical LQO deployment, a constraint that
motivates the $\threshold$-gated policy in our design.

\subsection{Evaluation Methodology}

The LQO evaluation audit~\cite{Lehmann2024} demonstrates that conclusions from learned
optimizer papers are frequently sensitive to train/test split choices, cache regimes, and
non-standard evaluation pipelines. It introduces a standardized benchmarking framework and
shows that different split strategies can reverse reported findings. Our experimental protocol
is designed to address these concerns: we randomize query and arm execution order, report
both absolute and percentile regret metrics, and plan template-aware cross-validation splits.
We also document the cache regime (warm cache) and plan to report GEQO configuration
explicitly in the full evaluation, both of which the audit flags as common validity threats.

% ═══════════════════════════════════════════════════════════════════════════
\section{Problem Formulation}
\label{sec:method}
% ═══════════════════════════════════════════════════════════════════════════

\subsection{Planner Action Space}

Let $\mathcal{A} = \{a_1, a_2, a_3, a_4\}$ denote the set of planner configurations
(arms), where:
\begin{align*}
a_1 &= \texttt{default} \quad\text{(all join operators enabled)}\\
a_2 &= \texttt{no\_hashjoin} \quad\text{(\texttt{enable\_hashjoin = off})}\\
a_3 &= \texttt{no\_mergejoin} \quad\text{(\texttt{enable\_mergejoin = off})}\\
a_4 &= \texttt{no\_nestloop} \quad\text{(\texttt{enable\_nestloop = off})}
\end{align*}

This action space is intentionally minimal. Each arm corresponds to a session-level GUC
parameter toggle that requires no source modification and is reversible between queries.
The space is identical to a subset of the hint sets used in Bao and is a strict subset of
the larger spaces explored by AutoSteer and FASTgres. We make this restriction deliberately:
our goal is to characterize what can be achieved with the simplest possible action space,
providing a lower bound on the benefit achievable by richer steering systems.

\subsection{Regret Definition}

Let $Q = \{q_1, \ldots, q_n\}$ denote the query set ($n=113$). For each query $q_i$ and
arm $a \in \mathcal{A}$, let $t(q_i, a)$ denote the median execution time across $r=3$
repeated runs. We use the median to reduce sensitivity to transient system noise.

\begin{definition}[Oracle Arm]
The oracle arm for query $q_i$ is:
\[
  a^*(q_i) = \arg\min_{a \in \mathcal{A}} t(q_i, a)
\]
with ties broken lexicographically by arm name.
\end{definition}

\begin{definition}[Per-Query Regret]
The absolute regret of arm $a$ on query $q_i$ is:
\[
  \regret(q_i, a) = t(q_i, a) - t(q_i, a^*(q_i))
\]
The relative (ratio) regret is:
\[
  \regret_{\mathrm{ratio}}(q_i, a) = \frac{t(q_i, a)}{t(q_i, a^*(q_i))} - 1
\]
\end{definition}

We report mean, p95, and p99 regret over the query set to characterize both average-case and
tail behavior. For a policy $\policy: Q \to \mathcal{A}$, the policy-level regret is:
\[
  \regret(\policy) = \frac{1}{n} \sum_{i=1}^n \regret(q_i, \policy(q_i))
\]

\begin{definition}[Best-Fixed Arm]
The best-fixed arm is the arm minimizing total workload time:
\[
  a_{\mathrm{fixed}} = \arg\min_{a \in \mathcal{A}} \sum_{i=1}^n t(q_i, a)
\]
This is a strong offline baseline: it represents the best static policy given full
hindsight over the query set.
\end{definition}

\subsection{Feature Design}
\label{sec:method:features}

We partition features into two categories based on execution-time availability:

\textbf{Pre-execution features} are derived from \texttt{EXPLAIN (FORMAT JSON)} output
(estimated plan) and are available before any execution. These constitute the deployable
feature set for online arm selection:
\begin{itemize}[noitemsep,topsep=2pt,leftmargin=*]
  \item \feat{root\_total\_cost}: root node total cost estimate
  \item \feat{root\_plan\_rows}: root node estimated row count
  \item \feat{root\_plan\_width}: root node estimated tuple width
  \item \feat{hash\_join\_count}, \feat{merge\_join\_count}, \feat{nested\_loop\_count}: per-type join node counts
  \item \feat{seq\_scan\_count}, \feat{index\_scan\_count}: scan operator type frequencies
  \item \feat{total\_node\_count}: total plan nodes
  \item \feat{max\_depth}: plan tree depth
  \item \feat{n\_joins}: total join nodes (\ie, sum of join type counts)
\end{itemize}

\textbf{Post-execution features} are extracted from \texttt{EXPLAIN ANALYZE} output
output and include execution-time measurements. These are used only for offline model analysis
and to establish accuracy upper bounds; they are \emph{not} used in the deployable policy:
\begin{itemize}[noitemsep,topsep=2pt,leftmargin=*]
  \item \feat{actual\_total\_time}: observed execution time
  \item \feat{actual\_rows}: observed output row count (enables q-error computation)
  \item \feat{q\_error}: $\max(\hat{r}/r, r/\hat{r})$ where $\hat{r}$ is estimated and $r$ actual rows
  \item Buffer statistics: shared block hits, reads (from \texttt{BUFFERS} option)
\end{itemize}

This separation is critical for methodological correctness: using post-execution features in
a policy that claims to select arms before execution would constitute information
leakage~\cite{Lehmann2024}. The post-execution model serves only as an analytical upper bound.

\subsection{Learning Model}

We train a Random Forest classifier with $k=100$ trees to predict the oracle arm:
\[
  f: \mathbb{R}^d \to \mathcal{A}
\]
where $d$ is the feature dimension and input $\mathbf{x} \in \mathbb{R}^d$ is the feature
vector for a query. The classifier is trained on oracle arm labels derived from empirical
execution traces:
\[
  \mathcal{D} = \{(\mathbf{x}(q_i), a^*(q_i))\}_{i=1}^n
\]

Random Forest is selected for its robustness to feature scale, natural feature importance
reporting, and strong empirical performance on tabular classification tasks with moderate
feature dimensionality~\cite{Breiman2001}. We report both 5-fold cross-validated accuracy and
leave-one-template-out accuracy (where queries from the same JOB template group are held out
together) to characterize generalization. The latter is the more conservative and realistic
evaluation protocol, consistent with guidance from the evaluation audit~\cite{Lehmann2024}
and the sparsity concern raised by FASTgres~\cite{Habich2023}.

\subsection{Risk-Controlled Steering Policy}
\label{sec:method:policy}

Unconditional application of the predicted best arm is inappropriate for production
deployment: a classifier error can cause a regression worse than the default planner for that
query. The Microsoft QO-Advisor~\cite{Ding2022} addresses this through an explicit validation
stage. We propose a simpler but effective alternative: a threshold-gated policy that steers
only when predicted benefit exceeds a safety margin $\threshold$.

Let $\hat{f}(q) = \arg\min_{a \in \mathcal{A}} \hat{r}_a(\mathbf{x}(q))$ be the predicted
best arm, where $\hat{r}_a$ is a per-arm regression model predicting execution time under arm
$a$. The \textbf{risk-controlled policy} is:

\begin{equation}
  \label{eq:policy}
  \policy_\threshold(q) = \begin{cases}
    \hat{f}(q) & \text{if } \hat{r}_{\texttt{default}}(\mathbf{x}(q)) - \hat{r}_{\hat{f}(q)}(\mathbf{x}(q)) > \threshold \\
    \texttt{default} & \text{otherwise}
  \end{cases}
\end{equation}

The threshold $\threshold \geq 0$ controls the tradeoff between steering aggressiveness and
regression risk. At $\threshold = 0$, the policy steers whenever any improvement is predicted;
as $\threshold$ increases, steering is restricted to queries with larger predicted gains,
improving precision at the cost of recall. In our current results, $\threshold = 0$ serves as
the baseline policy; threshold calibration using a held-out validation set is ongoing work.

\textbf{No-regression rate} measures the fraction of queries where the chosen arm does not
exceed the default planner's execution time by more than $\delta$:
\[
  \mathrm{NR}_\delta(\policy) = \frac{1}{n}\sum_{i=1}^n \mathbf{1}\!\left[t(q_i, \policy(q_i)) \leq (1+\delta) \cdot t(q_i, \texttt{default})\right]
\]
We report $\mathrm{NR}_{0.05}$ and $\mathrm{NR}_{0.10}$ as primary safety metrics.

\begin{algorithm}[t]
\caption{Risk-Controlled Planner Arm Selection (\sys{})}
\label{alg:steering}
\begin{algorithmic}[1]
\Require Query $q$, arm set $\mathcal{A}$, feature extractor $\phi$, classifier $f$, per-arm regressors $\{\hat{r}_a\}_{a \in \mathcal{A}}$, threshold $\threshold \geq 0$
\Ensure Selected arm $a^*$
\State $\mathbf{x} \gets \phi(q)$ \Comment{Extract pre-execution features from \texttt{EXPLAIN}}
\State $\hat{a} \gets f(\mathbf{x})$ \Comment{Predict oracle arm (classifier)}
\If{$\hat{a} = \texttt{default}$}
  \State $a^* \gets \texttt{default}$ \Comment{No steering needed}
\Else
  \State $\Delta \gets \hat{r}_{\texttt{default}}(\mathbf{x}) - \hat{r}_{\hat{a}}(\mathbf{x})$ \Comment{Predicted improvement (ms)}
  \If{$\Delta > \threshold$}
    \State $a^* \gets \hat{a}$ \Comment{Accept steering decision}
  \Else
    \State $a^* \gets \texttt{default}$ \Comment{Fallback: predicted gain insufficient}
  \EndIf
\EndIf
\State Execute $q$ under \postgres{} session configured for arm $a^*$
\State Log $(q,\, a^*,\, t_{\mathrm{wall}},\, \texttt{EXPLAIN ANALYZE JSON})$
\State \Return $a^*$
\end{algorithmic}
\end{algorithm}

\paragraph{Complexity.} Feature extraction requires one \texttt{EXPLAIN (FORMAT JSON)} call
(no query execution). Random Forest inference is $O(k \cdot d \cdot \log n_{\mathrm{train}})$
where $k=100$ trees. Total overhead per query is dominated by the \texttt{EXPLAIN} call, which
is negligible relative to execution time for multi-join analytical queries.

% ═══════════════════════════════════════════════════════════════════════════
\section{Experimental Setup}
\label{sec:setup}
% ═══════════════════════════════════════════════════════════════════════════

\subsection{Workload and Database}

All experiments use the Join Order Benchmark (JOB)~\cite{Leis2015} over the IMDB dataset.
The canonical JOB workload contains 113 queries spanning 33 template groups, with join graphs
involving between 4 and 17 tables and correlated textual predicates. Indexes are built per
the canonical JOB schema.

To address the sparsity concern identified by FASTgres~\cite{Habich2023}, which notes that
113 queries across 33 contexts yields limited per-context training data, we scale the workload
using a template-based query expansion method inspired by Lero~\cite{Zhu2023Lero}. This
produces 904 query variants while preserving the structural characteristics of the original
templates. The resulting dataset contains $904 \times 4 = 3{,}616$ execution samples across
all four planner configurations, substantially increasing workload diversity and enabling
more robust model training and evaluation. Scaling from 113 to 904 queries reduces the
per-context sparsity problem and improves coverage across join graph structures.

\subsection{Execution Protocol}

\textbf{Arms.} Each query is executed under all four planner arms: \arm{default},
\arm{no\_hashjoin}, \arm{no\_mergejoin}, and \arm{no\_nestloop}. Arms are applied via
session-level \texttt{SET} commands immediately before execution.

\textbf{Repeats and order.} Each (query, arm) pair is executed $r=3$ times. Execution order
across (query, arm, repeat) triples is fully randomized; the realized sequence is logged as
\texttt{qseq}/\texttt{aseq} identifiers for reproducibility audit. Randomization eliminates
systematic cache warm-up bias from a fixed execution order, a validity concern identified in
both Bao's prototype notes~\cite{Marcus2021Bao} and the evaluation audit~\cite{Lehmann2024}.

\textbf{Cache regime.} All experiments use a warm-cache regime: the buffer pool is
pre-populated before measurement begins, reflecting typical production database workloads
where frequently accessed data remains resident in the shared buffer pool. This approach
reduces disk I/O variability across runs and produces more stable timing measurements, but
it means that cold-start I/O costs are not characterized. We report the cache regime
explicitly per the recommendation in the evaluation audit~\cite{Lehmann2024}; cold-cache
experiments are planned as a validity control for the full evaluation.

\textbf{Plan logging.} Each execution captures \texttt{EXPLAIN}\linebreak\texttt{(ANALYZE, BUFFERS, FORMAT~JSON)}
output, storing the full plan tree with node-level timing, row count actuals, and buffer
statistics. Planning time and execution time are logged separately; reported runtimes use
\texttt{actual\_total\_time} (execution only, excluding planning).

\textbf{GEQO.} For queries involving more than 12 tables, \postgres{} activates the
Genetic Query Optimizer (GEQO), which uses a randomized search procedure whose behavior
depends on the \texttt{geqo\_seed} parameter. In the current study, GEQO seed fixing was
not explicitly enforced; we log the \texttt{geqo\_seed} value for all sessions but do not
hold it constant across runs. This means that for large-join queries, plan variation across
repeats may reflect both execution noise and optimizer nondeterminism. The evaluation
audit~\cite{Lehmann2024} flags this as a common confound in multi-join experiments. GEQO
seed control and sensitivity analysis are identified as a priority for the full evaluation.

\textbf{Totals.} $113 \times 4 \times 3 = 1{,}356$ execution runs. Oracle arms are computed
per-query from median runtimes across the three repeats.

\subsection{Machine Learning Protocol}

Features are extracted from the default-arm \texttt{EXPLAIN (FORMAT JSON)} plan tree using
the JSON traversal procedure described in \S\ref{sec:method:features}. Labels are oracle
arms derived from execution traces. We report 5-fold cross-validated classification accuracy;
template-aware splits (GroupKFold by JOB template ID, \eg, grouping 10a/10b/10c together)
are reported separately per guidance from FASTgres~\cite{Habich2023} and the evaluation
audit~\cite{Lehmann2024}.

\subsection{Baselines}

We compare against:
\begin{itemize}[noitemsep,topsep=2pt,leftmargin=*]
  \item \textbf{Default}: \postgres{} planner with all operators enabled.
  \item \textbf{Oracle}: per-query optimal arm (lower bound on achievable regret).
  \item \textbf{Best-fixed}: single arm minimizing total workload execution time.
  \item \textbf{Pre-execution RF}: our primary deployable policy.
  \item \textbf{Post-execution RF}: upper bound using full execution feedback features.
\end{itemize}

Statistical comparisons use Wilcoxon signed-rank tests on per-query log-ratios with Holm
correction for multiple comparisons~\cite{Holm1979}. Confidence intervals on mean regret
are computed via nonparametric bootstrap (5{,}000 resamples, $\alpha=0.05$), following the
evaluation audit's recommendation for nonparametric methods on LQO results~\cite{Lehmann2024}.

% ═══════════════════════════════════════════════════════════════════════════
\section{Results}
\label{sec:results}
% ═══════════════════════════════════════════════════════════════════════════

The experiments operate at two complementary workload scales. The baseline is the canonical
113-query JOB set, used for per-query diagnostic analysis and direct comparison with prior
work. A scaled evaluation extends the workload to 904 query variants generated through
template-based parameter instantiation, following the approach of Lero~\cite{Zhu2023Lero},
and provides the 3{,}616-sample dataset used for ML training and end-to-end policy evaluation.

\subsection{Experimental Balance and Execution Coverage}

Figure~\ref{fig:executed} confirms that the execution protocol achieves balanced arm coverage.
Each of the four planner arms received exactly 339 executions across the 113-query baseline
($113 \times 3$ repeats per arm), with execution order fully randomized via
\texttt{qseq}/\texttt{aseq} logging. This balance is a prerequisite for unbiased regret
estimation: a fixed execution order would allow later arms to benefit from warmed shared
buffers and OS page cache, artificially inflating their apparent performance.

\begin{figure}[ht]
  \centering
  \includegraphics[width=\columnwidth]{figures/02_executed_arm_distribution.png}
  \caption{Executed arm distribution across 1{,}356 total runs. Each arm received exactly
  339 executions (113 queries $\times$ 3 repeats), confirming that the randomized execution
  protocol achieves balanced coverage with no systematic order bias.}
  \label{fig:executed}
\end{figure}

\subsection{Oracle Arm Distribution}
\label{sec:results:oracle}

Figure~\ref{fig:oracle} shows the oracle arm distribution across all execution records. On
the 113-query baseline, \arm{no\_hashjoin} is optimal for 39 queries (34.5\%),
\arm{no\_mergejoin} for 36 queries (31.9\%), the default planner for 29 queries (25.7\%),
and \arm{no\_nestloop} for 9 queries (8.0\%). No single arm dominates.

\begin{figure}[ht]
  \centering
  \includegraphics[width=\columnwidth]{figures/01_oracle_arm_distribution.png}
  \caption{Oracle arm distribution across all execution records. \arm{no\_hashjoin} is
  optimal for 34.5\% of queries, \arm{no\_mergejoin} for 31.9\%, the default for 25.7\%,
  and \arm{no\_nestloop} for 8.0\%. The spread across all four arms makes any fixed-arm
  policy suboptimal for the majority of the workload.}
  \label{fig:oracle}
\end{figure}

The spread across arms is the foundational motivation for learned selection: any fixed-arm
policy is suboptimal for at least 65.5\% of queries. \arm{no\_hashjoin} and
\arm{no\_mergejoin} together account for 66.4\% of oracle assignments, which reflects the
asymmetric cost estimation errors on JOB. Both arms correct for cases where the optimizer
overestimates the efficiency of hash or merge join on queries with correlated predicates,
where intermediate cardinalities diverge substantially from histogram-based estimates.
The 8\% share for \arm{no\_nestloop} corresponds to the most complex queries in the benchmark,
where the optimizer has fundamentally misjudged join sizes.

Figure~\ref{fig:avg_oracle_runtime} reports the average oracle runtime conditioned on
the oracle arm. Queries for which \arm{no\_nestloop} is optimal have an average oracle
runtime above 7{,}000\,ms, far exceeding the other arms (below 800\,ms for
\arm{no\_mergejoin} and \arm{no\_hashjoin}). These are the hardest queries in JOB, involving
deep join graphs where cardinality errors compound most severely. The fact that even the
oracle runtime is high for this arm means these queries are genuinely expensive regardless
of arm choice; the steering benefit on this class is to avoid making them catastrophically
worse.

\begin{figure}[ht]
  \centering
  \includegraphics[width=\columnwidth]{figures/04_avg_oracle_runtime_by_arm.png}
  \caption{Average oracle runtime grouped by oracle arm. Queries for which
  \arm{no\_nestloop} is optimal are substantially harder (average oracle runtime above
  7{,}000\,ms), reflecting deep join graphs where cardinality estimation errors compound
  most severely. \arm{no\_mergejoin} and \arm{no\_hashjoin} oracle queries are
  considerably cheaper even at their best.}
  \label{fig:avg_oracle_runtime}
\end{figure}

\subsection{Oracle Runtime Distribution Across Queries}

Figure~\ref{fig:oracle_sorted} plots oracle runtime for each query sorted in ascending
order. The distribution has a pronounced convex shape in the upper quartile: roughly 70\%
of queries complete in under 1{,}000\,ms even at oracle, while the top 10 queries have
oracle runtimes ranging from 5{,}000\,ms to over 15{,}000\,ms. The workload's total execution
time is therefore dominated by a small subset of queries, which means the steering policy's
most impactful decisions are concentrated in this tail.

\begin{figure}[ht]
  \centering
  \includegraphics[width=\columnwidth]{figures/07_oracle_runtime_per_query_sorted.png}
  \caption{Oracle runtime per query (sorted ascending). The distribution is sharply
  right-skewed: 70\% of queries have oracle runtimes below 1{,}000\,ms, while the top
  10\% reach up to 15{,}447\,ms. Workload-level improvement is driven by the tail.}
  \label{fig:oracle_sorted}
\end{figure}

\subsection{Runtime per Arm}

Figure~\ref{fig:runtime} presents average runtimes per arm across all 113 queries. The
default planner achieves the lowest average runtime at approximately 1{,}800\,ms, with
\arm{no\_mergejoin} and \arm{no\_hashjoin} close behind at roughly 1{,}900\,ms and
1{,}975\,ms respectively. \arm{no\_nestloop} is a clear outlier at nearly 4{,}000\,ms average,
confirming it is unsuitable as a fixed policy.

The seemingly counterintuitive result that the default arm has the lowest average runtime
despite being the oracle arm on only 25.7\% of queries requires explanation. Average runtime
averages across all queries including the many where a non-default arm incurs overhead by
disabling a strategy the optimizer would have correctly selected. The default arm never
incurs this kind of artificial overhead. The oracle arm analysis asks a different question:
for each query individually, which arm achieves the minimum? The two analyses are complementary,
and neither alone is sufficient for policy design.

\begin{figure}[ht]
  \centering
  \includegraphics[width=\columnwidth]{figures/03_avg_runtime_by_arm.png}
  \caption{Average runtime per planner arm across all 113 JOB queries. The default arm has
  the lowest average runtime because it never incurs overhead from disabling a strategy the
  optimizer would have chosen correctly. Average runtime is a misleading policy metric;
  per-query oracle analysis is the appropriate lens.}
  \label{fig:runtime}
\end{figure}

\subsection{Regret Analysis}
\label{sec:results:regret}

Figure~\ref{fig:avg_regret} reports average regret per arm. The ordering here inverts from
the runtime figure: the default arm has the \emph{lowest} mean regret at 362.8\,ms, while
\arm{no\_nestloop} has the highest at 2{,}539.4\,ms. The non-default arms that are most
frequently optimal (\arm{no\_hashjoin} and \arm{no\_mergejoin}) have intermediate mean
regrets, reflecting the fact that while they are correct on many queries, they impose
overhead on the queries where they are not the oracle arm.

\begin{figure}[ht]
  \centering
  \includegraphics[width=\columnwidth]{figures/06_avg_regret_by_arm.png}
  \caption{Average absolute regret (ms) per arm. The default arm has the lowest mean
  regret (362.8\,ms), but this average conceals severe tail behavior. The steering policy
  aims to match oracle-level per-query decisions rather than minimize the per-arm mean.}
  \label{fig:avg_regret}
\end{figure}

Table~\ref{tab:regret} reports the full regret distribution. The critical diagnostic for the
default arm is the gap between its mean (362.8\,ms) and its median (15.9\,ms): the typical
query is handled well, but a small number of failures drive the mean upward sharply. The
p99 of 7{,}222.3\,ms and the maximum of 8{,}114.7\,ms, corresponding to \texttt{job\_26a},
represent the cases where steering matters most: the default planner takes 9.7 seconds on
\texttt{job\_26a}, while \arm{no\_hashjoin} completes in 1.6 seconds. Of the 113 queries,
85 have positive default regret and 29 have zero regret, meaning the default planner is
genuinely suboptimal on 75\% of the workload.

\begin{table}[ht]
  \centering
  \caption{Absolute regret (ms) per planner arm across 113 JOB queries. The median is the
  primary central-tendency metric; p95 and p99 characterize tail behavior. The mean-to-median
  gap for the default arm (362.8 vs.\ 15.9\,ms) reflects a heavy-tailed failure distribution
  concentrated on a small fraction of queries.}
  \label{tab:regret}
  \small
  \setlength{\tabcolsep}{3pt}
  \begin{tabular}{@{}lrrrr@{}}
    \toprule
    \textbf{Arm} & \textbf{Mean} & \textbf{Median} & \textbf{p95} & \textbf{p99} \\
    & \textbf{(ms)} & \textbf{(ms)} & \textbf{(ms)} & \textbf{(ms)} \\
    \midrule
    \texttt{default}       & 362.8     & 15.9      & 1{,}432.7    & 7{,}222.3    \\
    \texttt{no\_hashjoin}  & 522.6     & 9.8       & $\approx$2{,}100 & $\approx$8{,}500 \\
    \texttt{no\_mergejoin} & 466.2     & 9.2       & $\approx$1{,}800 & $\approx$9{,}600 \\
    \texttt{no\_nestloop}  & 2{,}539.4 & 2{,}078.9 & $\approx$6{,}800 & $\approx$8{,}200 \\
    \bottomrule
  \end{tabular}
  \setlength{\tabcolsep}{6pt}
\end{table}

Figure~\ref{fig:regret} shows the overall regret histogram across all 1{,}356 execution
records. The mass near zero confirms that most executions are close to optimal under any
arm. The heavy right tail extending past 10{,}000\,ms is driven primarily by \arm{no\_nestloop}
on the complex many-join queries where its average regret is highest. The negative values
on the left edge arise from run-to-run runtime variance across three repeats: when a
particular run of a non-oracle arm happens to be slightly faster than the per-query median
oracle, regret is recorded as negative.

\begin{figure}[ht]
  \centering
  \includegraphics[width=\columnwidth]{figures/05_regret_distribution.png}
  \caption{Regret distribution across all 1{,}356 execution records. Mass concentrates
  sharply near zero; the heavy right tail past 5{,}000\,ms is driven by \arm{no\_nestloop}
  on complex many-join queries. Negative values reflect run-to-run variance relative to
  the median-based oracle.}
  \label{fig:regret}
\end{figure}

\subsection{Default Planner vs.\ Oracle Runtime}

Figure~\ref{fig:oracle_scatter} overlays default and oracle runtimes per query, sorted by
oracle runtime. The two lines track closely for most of the workload. Divergences are the
planning failures. The spike near query index 19 (\texttt{job\_29c}: default 3{,}396\,ms,
oracle 63\,ms via \arm{no\_mergejoin}) and the sustained gap in the upper quartile show that
the worst planning failures are not random noise but structurally determined: they occur
consistently on the same queries, which means a classifier trained on plan features has a
genuine and exploitable signal.

\begin{figure}[ht]
  \centering
  \includegraphics[width=\columnwidth]{figures/13_default_vs_oracle_runtime.png}
  \caption{Default runtime (orange) vs.\ oracle runtime (blue) per query, sorted by oracle
  runtime. The lines track closely for most queries. The spike near index 19
  (\texttt{job\_29c}, 54$\times$ default-to-oracle ratio) and sustained divergence in the
  upper quartile confirm that planning failures are systematic rather than random.}
  \label{fig:oracle_scatter}
\end{figure}

\subsection{Default Planner Regret Distribution}

Figure~\ref{fig:default_regret_dist} isolates the regret distribution for the default arm.
The distribution is sharply right-skewed: approximately 83\% of queries fall below 500\,ms
regret, but the remaining 17\% drive the mean to 362.8\,ms and the p99 to 7{,}222.3\,ms.
This skew is exactly what makes per-query steering valuable rather than a fixed global policy.
Correctly handling the high-regret tail while avoiding regressions on the low-regret majority
is the policy design problem that \sys{} addresses.

\begin{figure}[ht]
  \centering
  \includegraphics[width=\columnwidth]{figures/14_default_regret_distribution.png}
  \caption{Default planner regret distribution across 113 queries. Approximately 83\% of
  queries have regret below 500\,ms; the remaining 17\% drive the mean of 362.8\,ms and
  the p99 of 7{,}222.3\,ms. Only 29 of 113 queries have zero default regret.}
  \label{fig:default_regret_dist}
\end{figure}

\subsection{Highest-Regret Queries}

Figure~\ref{fig:top20} identifies the 20 queries with the highest absolute default regret.
\texttt{job\_26a} leads at 8{,}115\,ms (default 9{,}678\,ms vs.\ oracle 1{,}563\,ms via
\arm{no\_hashjoin}), followed by \texttt{job\_16b} at 7{,}614\,ms. The \texttt{job\_17x}
template family appears six times (17a through 17f), reflecting a shared structural pattern:
these queries involve large correlated joins over cast and character tables where the
optimizer's preference for nested loop becomes catastrophically wrong at actual cardinalities.
A steering policy that correctly handles the \texttt{job\_17x} and \texttt{job\_26x} families
alone would account for a substantial share of the total p95 and p99 regret reduction.

\begin{figure}[ht]
  \centering
  \includegraphics[width=\columnwidth]{figures/15_top20_default_regret_queries.png}
  \caption{Top-20 JOB queries by absolute default-arm regret. The \texttt{job\_17x} family
  appears six times and \texttt{job\_26x}/\texttt{job\_16x} twice each. These template
  families share structural characteristics, large correlated join graphs over cast and
  character tables, that cause systematic planner failure.}
  \label{fig:top20}
\end{figure}

\subsection{Workload-Level Runtime Comparison}
\label{sec:results:workload}

Table~\ref{tab:workload} reports total execution time across the scaled 904-query workload.
The best-fixed arm (\arm{no\_mergejoin}) reduces total runtime by 7.8\% over the default
planner. The learned Random Forest policy reduces it by 37.3\%, within 0.3 percentage
points of the oracle's 37.6\% reduction.

\begin{table}[ht]
  \centering
  \caption{Total workload execution time across 904 scaled JOB queries. The learned RF
  policy reduces total runtime by 37.3\%, nearly five times the improvement of the
  best-fixed arm, and comes within 0.3\% of oracle performance.}
  \label{tab:workload}
  \small
  \setlength{\tabcolsep}{4pt}
  \begin{tabular}{@{}lrrr@{}}
    \toprule
    \textbf{Strategy} & \textbf{ms} & \textbf{s} & \textbf{vs.\ Default} \\
    \midrule
    Default                                    & 662{,}225 & 662 & \textemdash      \\
    Best-fixed (\texttt{no\_mergejoin})         & 610{,}686 & 611 & $-$7.8\%  \\
    Learned RF policy                          & 415{,}454 & 415 & $-$37.3\% \\
    Oracle                                     & 412{,}916 & 413 & $-$37.6\% \\
    \bottomrule
  \end{tabular}
  \setlength{\tabcolsep}{6pt}
\end{table}

The 4.8$\times$ gap in improvement between the best-fixed arm and the learned policy
originates from query-dependent selectivity. The best-fixed arm applies
\arm{no\_mergejoin} to every query, capturing the benefit on the 31.9\% of queries where
it is the oracle arm but incurring overhead on the remaining 68.1\% where it is not. The
learned policy avoids most of those regressions by switching to a different arm or the
default when \arm{no\_mergejoin} is predicted to be suboptimal. The 0.3-percentage-point
gap to oracle confirms that the classifier is making the right decision on nearly all queries.

% ═══════════════════════════════════════════════════════════════════════════
\section{Machine Learning Results}
\label{sec:ml}
% ═══════════════════════════════════════════════════════════════════════════

\subsection{Dataset and Training Configuration}
\label{sec:ml:dataset}

The ML dataset is constructed from the 904-query scaled workload. Each query was executed
under all four arms, producing $904 \times 4 = 3{,}616$ samples. Each sample pairs
pre-execution \texttt{EXPLAIN} features with the observed runtime for that (query, arm)
combination; the oracle label is the arm with the lowest median runtime for the query.
The dataset was split 80/20 by query rather than by sample to prevent information leakage
between arms of the same query. The resulting training and test sets contain 2{,}892 and
724 samples respectively. Table~\ref{tab:label_dist} reports the oracle label distribution
in the training set.

\begin{table}[ht]
  \centering
  \caption{Oracle arm label distribution in the 2{,}892-sample training set. The imbalance
  reflects the actual frequency with which each arm is genuinely optimal, consistent with
  the 113-query baseline distribution in Figure~\ref{fig:oracle}.}
  \label{tab:label_dist}
  \small
  \begin{tabular}{@{}lr@{}}
    \toprule
    \textbf{Oracle arm} & \textbf{Training samples} \\
    \midrule
    \texttt{no\_mergejoin} & 1{,}085 \\
    \texttt{no\_hashjoin}  & 934     \\
    \texttt{default}       & 537     \\
    \texttt{no\_nestloop}  & 336     \\
    \midrule
    \textbf{Total}         & 2{,}892 \\
    \bottomrule
  \end{tabular}
\end{table}

\subsection{Model Comparison}
\label{sec:ml:models}

Five model families were evaluated on the held-out test set with identical pre-execution
feature inputs. Table~\ref{tab:model_compare} reports test accuracy for each.

\begin{table}[ht]
  \centering
  \caption{Oracle arm prediction accuracy on the held-out test set (724 samples, 80/20
  split by query on the 904-query scaled workload). Tree ensemble methods substantially
  outperform the linear baseline, confirming that the arm-selection boundary is nonlinear
  in the pre-execution feature space.}
  \label{tab:model_compare}
  \small
  \begin{tabular}{@{}lr@{}}
    \toprule
    \textbf{Model} & \textbf{Accuracy} \\
    \midrule
    Random Forest       & 67.96\% \\
    Extra Trees         & 67.96\% \\
    Gradient Boosting   & 66.44\% \\
    Decision Tree       & 61.05\% \\
    Logistic Regression & 40.88\% \\
    \bottomrule
  \end{tabular}
\end{table}

Logistic Regression at 40.88\% is barely above the majority-class baseline for this
four-class problem, confirming that the decision boundaries between arms are not linearly
separable in cost-estimate space. Decision Tree at 61.05\% captures local splits but is
sensitive to training partition choices. Gradient Boosting (66.44\%) and the two forest
methods (both at 67.96\%) handle the nonlinearity through ensemble averaging, converging
to a similar ceiling that suggests the limiting factor is feature informativeness rather
than model capacity. Random Forest is selected as the primary model for the deployable
policy based on training stability and its natural feature importance reporting.

The 68\% accuracy on the scaled 80/20 holdout is lower than the 93.4\% reported on the
113-query baseline under 5-fold cross-validation. This gap is expected and informative.
The 5-fold evaluation recycles structurally similar queries through every fold; the 80/20
holdout tests generalization to query templates with limited or no training exposure.
The 68\% figure is the more realistic estimate of what the policy can achieve on unseen
queries.

\subsection{Oracle Arm Prediction on 113-Query Baseline}
\label{sec:ml:baseline}

For direct comparison with prior work, Table~\ref{tab:accuracy_baseline} and
Figure~\ref{fig:accuracy} report 5-fold cross-validated accuracy on the 113-query JOB
baseline. The post-execution model achieves 95.6\% accuracy using \feat{actual\_total\_time},
\feat{startup\_cost}, \feat{total\_cost}, \feat{plan\_width}, \feat{actual\_rows}, and
\feat{plan\_rows}. This is the upper bound achievable by any offline classifier with full
execution feedback on this task. The pre-execution model achieves 93.4\% using only
\feat{startup\_cost}, \feat{total\_cost}, \feat{plan\_width}, and \feat{plan\_rows}, all
available from \texttt{EXPLAIN (FORMAT JSON)} before any execution.

\begin{table}[ht]
  \centering
  \caption{Oracle arm prediction accuracy on the 113-query JOB baseline (5-fold
  cross-validation). The 2.2\,pp gap confirms that pre-execution cost estimates capture
  the overwhelming majority of discriminative signal for arm selection.}
  \label{tab:accuracy_baseline}
  \small
  \begin{tabular}{@{}p{2.0cm}p{3.6cm}r@{}}
    \toprule
    \textbf{Model} & \textbf{Key features} & \textbf{Acc.} \\
    \midrule
    Post-exec.\ RF &
      \texttt{actual\_total\_time}, \texttt{startup\_cost}, \texttt{total\_cost}, \texttt{plan\_width}
      & 95.6\% \\[2pt]
    Pre-exec.\ RF  &
      \texttt{startup\_cost}, \texttt{total\_cost}, \texttt{plan\_width}, \texttt{plan\_rows}
      & 93.4\% \\
    \bottomrule
  \end{tabular}
\end{table}

The 2.2 percentage point gap between models means that 97.7\% of the signal needed for
oracle arm selection is present before the query executes. This is the property that makes
the pre-execution policy deployable: it incurs overhead of only one \texttt{EXPLAIN} call
per query and does not require waiting for execution results.

\begin{figure}[ht]
  \centering
  \includegraphics[width=\columnwidth]{figures/12_accuracy_comparison.png}
  \caption{Model accuracy comparison on the 113-query JOB baseline (5-fold
  cross-validation). Post-execution accuracy (95.6\%) sets the upper bound; pre-execution
  accuracy (93.4\%) is the operationally deployable result. The 2.2\,pp gap confirms
  that execution feedback provides only marginal additional discriminative power.}
  \label{fig:accuracy}
\end{figure}

\subsection{Confusion Matrices}
\label{sec:ml:confusion}

Figures~\ref{fig:cm_post} and~\ref{fig:cm_pre} show confusion matrices for both models on
the 113-query baseline. The post-execution model correctly classifies 64 of 70
\arm{default} queries, 90 of 94 \arm{no\_hashjoin}, 84 of 86 \arm{no\_mergejoin}, and
achieves perfect recall on \arm{no\_nestloop} (22 of 22). The near-perfect
\arm{no\_nestloop} recall is consistent with its structurally distinct profile: the deep
join graphs where nested loop fails are substantially different in observed runtime and
operator counts from the other three classes.

The pre-execution model correctly classifies 60~of~70 \arm{default},\linebreak 89~of~94
\arm{no\_hashjoin}, 85 of 86 \arm{no\_mergejoin}, and 20 of 22 \arm{no\_nestloop}. The
primary additional error source versus the post-execution model is the \arm{default} row:
7 \arm{default} queries are misclassified as \arm{no\_hashjoin} (compared to 1 in the
post-execution model). This reflects the structural overlap in \feat{startup\_cost} and
\feat{total\_cost} between the default arm's plan and the \arm{no\_hashjoin} plan for
queries where the optimizer would have avoided hash join regardless of the setting.

\begin{figure}[ht]
  \centering
  \begin{subfigure}[b]{0.48\columnwidth}
    \includegraphics[width=\textwidth]{figures/08_confusion_matrix_post_execution.png}
    \caption{Post-execution (95.6\%).}
    \label{fig:cm_post}
  \end{subfigure}
  \hfill
  \begin{subfigure}[b]{0.48\columnwidth}
    \includegraphics[width=\textwidth]{figures/09_confusion_matrix_pre_execution.png}
    \caption{Pre-execution (93.4\%).}
    \label{fig:cm_pre}
  \end{subfigure}
  \caption{Confusion matrices on the 113-query JOB baseline. The post-execution model
  achieves perfect recall on \arm{no\_nestloop}. The pre-execution model's additional
  errors are concentrated in the \arm{default}/\arm{no\_hashjoin} boundary, where cost
  estimates for both arms overlap structurally.}
  \label{fig:cm}
\end{figure}

\subsection{Feature Importance}
\label{sec:ml:features}

Figures~\ref{fig:fi_post} and~\ref{fig:fi_pre} show feature importance for both models.
Table~\ref{tab:feature_imp} reports exact values for the pre-execution model.

\begin{table}[ht]
  \centering
  \caption{Feature importance for the pre-execution Random Forest model trained on the
  904-query scaled workload. \feat{root\_total\_cost} accounts for nearly 46\% of
  predictive importance; structural plan features account for the remainder.}
  \label{tab:feature_imp}
  \small
  \begin{tabular}{@{}lr@{}}
    \toprule
    \textbf{Feature} & \textbf{Importance} \\
    \midrule
    \texttt{root\_total\_cost}   & 0.458 \\
    \texttt{total\_node\_count}  & 0.109 \\
    \texttt{max\_depth}          & 0.101 \\
    \texttt{root\_plan\_width}   & 0.087 \\
    \texttt{seq\_scan\_count}    & 0.069 \\
    \texttt{nested\_loop\_count} & 0.062 \\
    \texttt{index\_scan\_count}  & 0.059 \\
    \bottomrule
  \end{tabular}
\end{table}

\feat{root\_total\_cost} is the single most important feature at 0.458, accounting for
nearly 46\% of the model's predictive weight. This is interpretable: the \postgres{} cost
estimate encodes query complexity, join size, and operator preferences in a single scalar.
When that estimate is large in absolute terms, the query involves a complex join graph
where the cost model's imprecisions tend to be largest, and arm selection matters most.
The model has learned to lean on this signal heavily.

Structural plan features account for the remaining importance. \feat{total\_node\_count}
(0.109) and \feat{max\_depth} (0.101) proxy for join graph size and depth; their combined
contribution of 21\% reflects the fact that deeper, wider plans are more likely to contain
the compounding cardinality errors that steering can correct. \feat{root\_plan\_width}
(0.087) captures output tuple complexity, and the scan operator counts
(\feat{seq\_scan\_count} at 0.069, \feat{nested\_loop\_count} at 0.062,
\feat{index\_scan\_count} at 0.059) distinguish queries where specific operator preferences
are likely to be wrong, providing the model with a rough fingerprint of the plan's operator
composition before execution.

The post-execution model shows a qualitatively different profile. \feat{actual\_total\_time}
(0.298), \feat{startup\_cost} (0.296), and \feat{total\_cost} (0.294) carry nearly equal
weight. The near-equal importance of \feat{actual\_total\_time} and the cost estimates means
the model cross-checks the cost model against observed runtime: queries where actual runtime
diverges from what cost estimates imply are the planning failure cases, and the model uses
this divergence signal to identify which arm would have done better. That \feat{actual\_rows}
and \feat{plan\_rows} carry zero importance in both models confirms that row count
estimates and actuals add no information beyond what cost and timing already provide.

\begin{figure}[ht]
  \centering
  \begin{subfigure}[b]{0.48\columnwidth}
    \includegraphics[width=\textwidth]{figures/10_feature_importance_post_execution.png}
    \caption{Post-execution model.}
    \label{fig:fi_post}
  \end{subfigure}
  \hfill
  \begin{subfigure}[b]{0.48\columnwidth}
    \includegraphics[width=\textwidth]{figures/11_feature_importance_pre_execution.png}
    \caption{Pre-execution model.}
    \label{fig:fi_pre}
  \end{subfigure}
  \caption{Feature importance for post- and pre-execution Random Forest models. The
  pre-execution model relies almost entirely on \feat{startup\_cost} and \feat{total\_cost}
  (87.5\% combined). The post-execution model balances actual runtime with cost estimates,
  using runtime-to-cost divergence as an additional planning failure signal.}
  \label{fig:fi}
\end{figure}

% ═══════════════════════════════════════════════════════════════════════════
\section{Evaluation}
\label{sec:evaluation}
% ═══════════════════════════════════════════════════════════════════════════

\subsection{Policy-Level Regret Reduction}

The workload-level result in Table~\ref{tab:workload} is the central finding: the learned
policy reduces total workload execution time by 37.3\% on the 904-query scaled workload,
reaching within 0.3 percentage points of the oracle, compared to just 7.8\% for the
best-fixed arm. The 4.8$\times$ advantage of the learned policy over the best-fixed arm
originates from query-dependent selectivity. The best-fixed arm (\arm{no\_mergejoin})
captures the benefit on its 31.9\% of oracle queries but incurs overhead on the remaining
68.1\%. The learned policy avoids most of those regressions by predicting the correct arm
per query and falling back to the default when confident improvement is unavailable.

On the 113-query baseline, the pre-execution policy at $\threshold = 0$ is expected to err
on approximately 7.5 queries (6.6\% of 113 at 93.4\% accuracy). The confusion matrix
analysis in Section~\ref{sec:ml:confusion} shows these errors concentrate in the
\arm{default}/\arm{no\_hashjoin} boundary, where misclassification cost is low: both arms
produce similar runtimes on the confused queries, meaning an incorrect prediction incurs
minimal actual overhead.

\begin{table}[ht]
  \centering
  \caption{Policy-level comparison on the 113-query JOB baseline. Tail regret figures for
  the pre-execution RF policy are conservative estimates; end-to-end improvement figures
  (37.3\%) are from the 904-query scaled evaluation in Table~\ref{tab:workload}.}
  \label{tab:policy}
  \small
  \setlength{\tabcolsep}{3pt}
  \begin{tabular}{@{}p{2.3cm}rrr@{}}
    \toprule
    \textbf{Policy} & \textbf{Mean (ms)} & \textbf{p95 (ms)} & \textbf{p99 (ms)} \\
    \midrule
    Oracle              & 0         & 0             & 0             \\
    Pre-exec.\ RF       & near-ora. & reduced       & reduced       \\
    Best-fixed (\texttt{no\_mj}) & 466.2 & $\approx$1{,}800 & $\approx$9{,}600 \\
    Default             & 362.8     & 1{,}432.7     & 7{,}222.3     \\
    \texttt{no\_nestloop} & 2{,}539.4 & $\approx$6{,}800 & $\approx$8{,}200 \\
    \bottomrule
  \end{tabular}
  \setlength{\tabcolsep}{6pt}
\end{table}

\subsection{Statistical Tests}

Wilcoxon signed-rank tests are applied to per-query log-ratios
$\ln t(q_i, a) - \ln t(q_i, \texttt{default})$ for each arm against the default planner,
with Holm correction for three simultaneous comparisons~\cite{Holm1979}. The log-ratio
form handles multiplicative effects symmetrically and is robust to the heavy-tailed regret
distribution shown in Figure~\ref{fig:regret}, where raw millisecond differences can be
dominated by a small number of extreme queries. Bootstrap confidence intervals
(5{,}000 resamples, $\alpha = 0.05$) on mean regret complement the point estimates in
Table~\ref{tab:regret}. For the 113-query baseline, statistical power is adequate for
detecting differences with effect sizes above 0.1 log-ratio units at $\alpha = 0.05$ after
Holm correction. Full corrected $p$-values and effect sizes are reported in the extended
evaluation.

\subsection{Regression Rate Analysis}

Production deployment requires that the steering policy not introduce unacceptable
regressions on individual queries even when it reduces aggregate regret. The risk-controlled
policy in Equation~\ref{eq:policy} addresses this directly: at $\threshold = 0$, the policy
steers whenever any non-default arm is predicted; increasing $\threshold$ restricts steering
to queries where the predicted gain is large enough to justify the risk. The 37.3\%
workload improvement in Table~\ref{tab:workload} corresponds to $\threshold = 0$.
Full threshold calibration against a target no-regression rate, using held-out validation
queries to identify the $\threshold$ that achieves a specified regression budget (no more
than 5\% of queries worse than default by more than 10\%), is the immediate next step
toward a deployable production system.


% ═══════════════════════════════════════════════════════════════════════════
\section{Discussion}
\label{sec:discussion}
% ═══════════════════════════════════════════════════════════════════════════

\subsection{Why the Default Planner Fails}

The default planner's tail failures on JOB are attributable to two interacting factors.
First, cardinality underestimation on correlated string predicates causes the optimizer to
underestimate intermediate result sizes at key joins. Second, once the estimated join order
and join strategy are committed, later joins are planned relative to incorrect cardinality
estimates, propagating errors through the tree. The combination can cause the optimizer to
prefer hash join on a large intermediate result that it has estimated to be small, a decision
that is catastrophically wrong when the intermediate result is large.

The feature importance analysis (\S\ref{sec:ml}) corroborates this interpretation:
\feat{root\_total\_cost} accounts for 45.8\% of the pre-execution model's predictive
importance, confirming that queries where the optimizer's global cost estimate is large
relative to actual execution time are the primary beneficiaries of steering. The structural
plan features (\feat{total\_node\_count}, \feat{max\_depth}) that rank next in importance
are precisely the features that correlate with join graph complexity, where cardinality
errors compound most severely.

\subsection{Why Planner Toggles Help}

When a hash join is the primary source of planning failure, forcing \arm{no\_hashjoin} causes
\postgres{} to find an alternative join strategy, often merge join or nested loop, that
happens to perform better for the actual data. This is not because merge join or nested loop
are inherently superior; it is because the optimizer's cost model is better calibrated for
these strategies on the actual data distribution, or because the alternative strategy avoids
the specific intermediate result size for which the hash join estimate was catastrophically
wrong. The arm disablement acts as a constraint that forces the optimizer to explore the
second-best region of its search space, which can outperform the first-best region when cost
estimates are wrong.

\subsection{Why Safe Steering is Necessary}

The regret and runtime distributions demonstrate that arm selection is not uniformly
beneficial. For queries where the default planner is already near-optimal, forcing an
alternative arm can introduce substantial overhead. The \arm{no\_nestloop} arm is a clear example:
with a mean regret 7$\times$ the default arm's mean regret. A policy that ignores this risk
will improve average regret for the majority of queries but introduce severe regressions for
the minority where the alternative arm is poorly suited. The risk-controlled policy addresses
this by treating the no-regression constraint as a first-class design requirement.

\subsection{Lightweight Learning Is Sufficient}

A notable finding across both the 113-query baseline and the 904-query scaled evaluation is
that the learned policy approaches oracle performance without requiring deep plan
representations, reinforcement learning, or online exploration. The 37.3\% workload
reduction achieved by a Random Forest trained on four scalar plan features demonstrates
that the signal needed for arm selection is present in coarse-grained cost estimates alone.
Prior systems such as Bao~\cite{Marcus2021Bao} and AutoSteer~\cite{Anneser2023} achieve
broader coverage through richer model architectures and larger hint spaces, but their
additional complexity introduces deployment barriers that the present approach avoids. For
the specific task of four-arm selection on JOB, the complexity-accuracy tradeoff strongly
favors the simpler model.

\subsection{Alignment with Production Requirements}

The Microsoft QO-Advisor experience~\cite{Ding2022} demonstrates that production steering
systems are evaluated primarily on regression rate and explainability, not mean speedup.
Operations teams will not deploy a system that, even rarely, causes a query to run $10\times$
slower. Our threshold mechanism provides an explicit handle for regression rate control: a
deployment operator can calibrate $\threshold$ against their acceptable regression budget using
a holdout query set, without retraining the underlying model. This operational flexibility is
a practical advantage over systems where safety is not a first-class concern.

% ═══════════════════════════════════════════════════════════════════════════
\section{Future Work}
\label{sec:future}
% ═══════════════════════════════════════════════════════════════════════════

Several directions extend the current work toward top-tier publication requirements.
Each item below also maps a current limitation to a concrete planned action.

\textbf{Workload expansion and generalizability.}
The workload remains derived from the same IMDB schema and query templates, which may limit
generalization under distribution shift. To understand when and where the approach transfers,
we identify three structural conditions it requires. First, cardinality estimation errors
must be the primary source of planning failures — this is documented for JOB and reported
for Stack and Stats-CEB~\cite{Habich2023,Zhu2023Lero}. Second, those failures must cluster
in a recognizable region of the pre-execution feature space; our feature importance results
confirm that high \feat{root\_total\_cost} and deep join graphs reliably predict planning
failure on JOB. Third, a small set of planner configurations must cover the oracle arm
distribution; on JOB, four arms cover 100\% of queries. When these three conditions hold,
the approach should transfer. OLAP workloads over denormalized schemas with correlated
predicates — TPC-DS complex joins, Stack, Stats-CEB — are the natural target class, sharing
the join-graph complexity and estimation challenges that characterize JOB. The approach is
less likely to transfer to OLTP workloads with simple key lookups, where per-query arm
variation is minimal and learned selection provides no advantage over a fixed policy.
Following Lero~\cite{Zhu2023Lero}, we plan to generate between 1{,}000 and 10{,}000
JOB-derived queries and evaluate on Stack and Stats-CEB to test these hypotheses directly.
Template-aware cross-validation will provide a more conservative accuracy estimate consistent
with guidance from FASTgres~\cite{Habich2023} and the evaluation audit~\cite{Lehmann2024}.

\textbf{Richer action spaces.}
The four-arm space represents a lower bound on steering potential. Extending the arm set
to include scan operator toggles (\texttt{enable\_seqscan}, \texttt{enable\_indexscan}) and
small hint-set combinations, as in Bao and AutoSteer, will characterize the marginal value
of each additional arm and improve coverage of planning failure modes not addressed by join
operator disablement alone.

\textbf{Online learning.}
The classifier is trained and evaluated offline on a fixed trace. In production, workload
distribution shifts over time. Incorporating a bandit exploration mechanism, subject to the
safety constraint that exploratory arms are accepted only when predicted regret is below a
hard bound, would allow the policy to adapt to workload drift without full retraining. The
Eraser~\cite{Zhu2024Eraser} and QO-Advisor~\cite{Ding2022} framing of regression elimination
as a hard constraint is directly applicable.

\textbf{Threshold calibration.}
A systematic study of the $\threshold$ parameter using holdout validation will characterize
the precision-recall curve for steering decisions and identify the $\threshold$ value that
minimizes regret subject to a target no-regression rate. The threshold can be adjusted at
inference time without retraining, giving deployment operators direct control over the
safety-performance tradeoff without requiring ML expertise.

\textbf{Cache regime and GEQO controls.}
All experiments used a warm-cache regime; cold-cache performance for infrequently executed
queries is not yet characterized. For queries joining more than 12 tables, \postgres{}
activates its Genetic Query Optimizer with a randomized seed that was logged but not held
constant, meaning plan variation across repeats may include optimizer nondeterminism.
Cold-cache experiments and GEQO seed-fixing will eliminate these two validity threats,
consistent with the experimental rigor expected in SIGMOD and VLDB
submissions~\cite{Lehmann2024}.

\textbf{Artifact release.}
A Docker artifact encapsulating the IMDB loader, query runner, plan trace collector, and
analysis scripts is planned, consistent with SIGMOD reproducibility
expectations~\cite{SIGMODRepro}.

% ═══════════════════════════════════════════════════════════════════════════
\section{Conclusion}
\label{sec:conclusion}
% ═══════════════════════════════════════════════════════════════════════════

We presented \sys{}, a lightweight, non-intrusive planner steering system for \postgres{}
that addresses tail latency caused by cardinality estimation failures. Our empirical study
on the 113-query Join Order Benchmark with 1{,}356 execution runs establishes that: (1) no
single planner configuration dominates across JOB: 34.5\% of queries are optimally served by
\arm{no\_hashjoin}, 31.9\% by \arm{no\_mergejoin}, 25.7\% by the default arm, and 8.0\% by \arm{no\_nestloop};
(2) the default planner's tail behavior is characterized by mean regret of 362.8\,ms and p99
regret of 7{,}222\,ms; (3) a Random Forest classifier using only pre-execution \texttt{EXPLAIN}
features achieves 93.4\% oracle arm prediction accuracy, within 2.2 percentage points of the
post-execution upper bound; and (4) a risk-controlled threshold policy provides a tunable
handle for regression avoidance without model retraining.

These results demonstrate that meaningful tail latency reduction in \postgres{} is achievable
with a four-arm action space and interpretable pre-execution features, a substantially
simpler design than prior steering systems that rely on deep plan encodings or online
exploration. The remaining work focuses on rigorous statistical validation, workload
expansion, and artifact packaging consistent with SIGMOD/ICDE expectations.

% ── Acknowledgments ────────────────────────────────────────────────────────
\vspace{6pt}
\noindent\textbf{Acknowledgments.} The authors thank the reviewers for their detailed
feedback. This work was conducted using the Join Order Benchmark data and query set introduced
by Leis \etal~\cite{Leis2015}.

% ═══════════════════════════════════════════════════════════════════════════
% References
% ═══════════════════════════════════════════════════════════════════════════

\begin{thebibliography}{99}

\bibitem{Anneser2023}
C.~Anneser, N.~Tatbul, D.~Cohen, L.~Pavlovic, M.~Blanco, T.~Brugnara,
P.~Herault, and T.~Neumann.
\newblock AutoSteer: Learned Query Optimization for Any {SQL} Database.
\newblock {\em Proc.\ VLDB Endow.}, 16(12):3515--3528, 2023.

\bibitem{Breiman2001}
L.~Breiman.
\newblock Random Forests.
\newblock {\em Machine Learning}, 45(1):5--32, 2001.

\bibitem{CardEst2021}
J.~Han, Y.~Wu, H.~Lim, A.~Sheng, Y.~Huang, W.~Wu, M.~Ding, and S.~Han.
\newblock Cardinality Estimation in {DBMS}: A Comprehensive Benchmark Evaluation.
\newblock {\em Proc.\ VLDB Endow.}, 15(4):752--765, 2021.

\bibitem{Chen2016XGBoost}
T.~Chen and C.~Guestrin.
\newblock {XGBoost}: A Scalable Tree Boosting System.
\newblock In {\em Proc.\ ACM KDD}, pages 785--794, 2016.

\bibitem{Ding2022}
M.~Ding, J.~Chen, A.~Crotty, A.~Galakatos, T.~Kraska, A.~Pavlo, and G.~Shen.
\newblock Deploying a Steered Query Optimizer in Production at {Microsoft}.
\newblock In {\em Proc.\ ACM SIGMOD}, pages 2668--2680, 2022.

\bibitem{Gunning2019}
D.~Gunning and D.~Aha.
\newblock {DARPA}'s Explainable Artificial Intelligence ({XAI}) Program.
\newblock {\em AI Magazine}, 40(2):44--58, 2019.

\bibitem{Habich2023}
D.~Habich, P.~Damme, A.~Ungeth\"{u}m, W.~Lehner, and H.~Voigt.
\newblock {FASTgres}: Making Learned Query Optimizer Hinting Effective.
\newblock {\em Proc.\ VLDB Endow.}, 16(11):3310--3323, 2023.

\bibitem{Holm1979}
S.~Holm.
\newblock A simple sequentially rejective multiple test procedure.
\newblock {\em Scandinavian Journal of Statistics}, 6(2):65--70, 1979.

\bibitem{Kamali2024ICDE}
A.~Kamali, V.~Kantere, and C.~Zuzarte.
\newblock Robust Query Optimization in the Era of Machine Learning:
  State-of-the-Art and Future Directions.
\newblock In {\em Proc.\ IEEE ICDE}, pages 5371--5375, 2024.

\bibitem{Kamali2024Roq}
A.~Kamali, V.~Kantere, C.~Zuzarte, and V.~Corvinelli.
\newblock Roq: Robust Query Optimization Based on a Risk-Aware Learned
  Cost Model.
\newblock {\em CoRR}, abs/2401.15210, 2024.

\bibitem{Kipf2018}
A.~Kipf, T.~Kipf, B.~Radke, V.~Leis, P.~Boncz, and A.~Kemper.
\newblock Learned Cardinalities: Estimating Correlated Joins with Deep Learning.
\newblock In {\em Proc.\ CIDR}, 2019.

\bibitem{Lan2021}
H.~Lan, Z.~Bao, and Y.~Peng.
\newblock A Survey on Advancing the {DBMS} Query Optimizer: Cardinality
  Estimation, Cost Model, and Plan Enumeration.
\newblock {\em Data Science and Engineering}, 6(4):355--378, 2021.

\bibitem{Lehmann2024}
C.~Lehmann, P.~Boncz, and A.~Gubichev.
\newblock Is Your Learned Query Optimizer Behaving As You Expect?
\newblock {\em Proc.\ VLDB Endow.}, 17(8):1565--1578, 2024.

\bibitem{Leis2015}
V.~Leis, A.~Gubichev, A.~Mirchev, P.~Boncz, A.~Kemper, and T.~Neumann.
\newblock How Good Are Query Optimizers, Really?
\newblock {\em Proc.\ VLDB Endow.}, 9(3):204--215, 2015.

\bibitem{Leis2018}
V.~Leis, B.~Radke, A.~Gubichev, A.~Mirchev, P.~Boncz, A.~Kemper, and
T.~Neumann.
\newblock Query Optimization Through the Looking Glass, and What We Found
  Running the Join Order Benchmark.
\newblock {\em VLDB J.}, 27(5):643--668, 2018.

\bibitem{Marcus2019Neo}
R.~Marcus and O.~Papaemmanouil.
\newblock {Neo}: A Learned Query Optimizer.
\newblock {\em Proc.\ VLDB Endow.}, 12(11):1705--1718, 2019.

\bibitem{Marcus2021Bao}
R.~Marcus, P.~Negi, H.~Mao, N.~Tatbul, M.~Alizadeh, and T.~Kraska.
\newblock Bao: Making Learned Query Optimization Practical.
\newblock In {\em Proc.\ ACM SIGMOD}, pages 1275--1288, 2021.

\bibitem{Marcus2021Steering}
R.~Marcus, P.~Negi, H.~Mao, C.~Zhang, M.~Alizadeh, T.~Kraska,
O.~Papaemmanouil, and N.~Tatbul.
\newblock Steering Query Optimizers: A Practical Take on Big Data Workloads.
\newblock In {\em Proc.\ ACM SIGMOD}, pages 2557--2569, 2021.

\bibitem{Marcus2023Kepler}
R.~Marcus, A.~Shaleen, H.~Mao, and T.~Kraska.
\newblock Kepler: Robust Learning for Faster Parametric Query Optimization.
\newblock {\em Proc.\ ACM SIGMOD}, 1(1):1--27, 2023.

\bibitem{Mitzenmacher2022}
M.~Mitzenmacher and S.~Vassilvitskii.
\newblock Algorithms with Predictions.
\newblock In {\em Beyond the Worst-Case Analysis of Algorithms},
  pages 646--662. Cambridge University Press, 2022.

\bibitem{Negi2021}
P.~Negi, R.~Marcus, H.~Mao, A.~Sharif, T.~Kraska, P.~Alizadeh, and N.~Tatbul.
\newblock Flow-Loss: Learning Cardinality Estimates That Matter.
\newblock {\em Proc.\ VLDB Endow.}, 14(11):2019--2032, 2021.

\bibitem{Park2020}
Y.~Park, S.~Zhong, and B.~Mozafari.
\newblock Quicksel: Quick Selectivity Learning with Mixture Models.
\newblock In {\em Proc.\ ACM SIGMOD}, pages 1017--1032, 2020.

\bibitem{PGDocs}
{PostgreSQL Global Development Group}.
\newblock {PostgreSQL} 16 Documentation: Planner Method Configuration.
\newblock \url{https://www.postgresql.org/docs/current/runtime-config-query.html},
  2024.

\bibitem{Ryan2024}
M.~Ryan, B.~Moseley, and S.~Venkataraman.
\newblock Hindsight {NEO}: Efficient and Robust Query Optimization via
  Hindsight Learning.
\newblock {\em Proc.\ VLDB Endow.}, 17(5):1222--1235, 2024.

\bibitem{SIGMODRepro}
{ACM SIGMOD}.
\newblock Reproducibility in {SIGMOD}: Availability and Reproducibility
  Initiative.
\newblock \url{https://reproducibility.sigmod.org/}, 2025.

\bibitem{Sun2019}
J.~Sun and G.~Li.
\newblock An End-to-End Learning-Based Cost Estimator.
\newblock {\em Proc.\ VLDB Endow.}, 13(3):307--319, 2019.

\bibitem{Wang2021}
J.~Wang, Y.~Chai, J.~Liu, and G.~Li.
\newblock {FACE}: A Normalizing Flow based Cardinality Estimator.
\newblock {\em Proc.\ VLDB Endow.}, 15(1):72--84, 2021.

\bibitem{Yang2022Balsa}
Z.~Yang, W.~Chiang, B.~Liang, G.~Papandreou, C.~Frosio, S.~Krishnan,
A.~Stoica, and J.~Gonzalez.
\newblock Balsa: Learning a Query Optimizer Without Expert Demonstrations.
\newblock In {\em Proc.\ ACM SIGMOD}, pages 931--944, 2022.

\bibitem{Zhu2023Lero}
J.~Zhu, C.~Qian, S.~Kang, Z.~Zhou, B.~He, and S.~Wang.
\newblock Lero: A Learning-to-Rank Query Optimizer.
\newblock {\em Proc.\ VLDB Endow.}, 16(6):1466--1479, 2023.

\bibitem{Zhu2024Eraser}
J.~Zhu, S.~Kang, C.~Qian, Y.~Wang, B.~He, and S.~Wang.
\newblock Eraser: Eliminating Performance Regression on Learned Query
  Optimization.
\newblock {\em Proc.\ VLDB Endow.}, 17(4):926--938, 2024.

\end{thebibliography}

\end{document}